\section{Throughput Optimality}
\label{sec:throughput}

We prove that BPE achieves throughput-optimal payment allocation within the
capacity region, following the Lyapunov drift framework of
Neely~\cite{neely2010stochastic}.

\subsection{Capacity Region}

\begin{definition}[Capacity Region]
The capacity region $\Lambda^*$ is the set of arrival rate vectors
$\boldsymbol{\lambda} = (\lambda_\tasktype)_{\tasktype \in \mathcal{T}}$ such that
there exists a feasible allocation $\{F(k, \tasktype)\}$ satisfying:
\begin{align}
\sum_{k \in K_\tasktype} F(k, \tasktype) &= \lambda_\tasktype &&\forall \tasktype \in \mathcal{T} \\
F(k, \tasktype) &\leq \Csmooth(k, \tasktype) &&\forall k, \tasktype \\
F(k, \tasktype) &\geq 0 &&\forall k, \tasktype
\end{align}
\end{definition}

\subsection{Lyapunov Function}

Define the Lyapunov function as the sum of squared virtual queue backlogs:
\begin{equation}
L(t) = \sum_{\tasktype \in \mathcal{T}} \sum_{k \in K_\tasktype} Q(k, t, \tasktype)^2
\end{equation}

The one-step conditional Lyapunov drift is:
\begin{equation}
\Delta(t) = \E\bigl[L(t+1) - L(t) \;\big|\; \mathbf{Q}(t)\bigr]
\end{equation}

\subsection{Main Result}

\begin{theorem}[Throughput Optimality of BPE]
\label{thm:main}
For any arrival rate vector $\boldsymbol{\lambda}$ strictly inside the capacity
region $\Lambda^*$ (i.e., $\boldsymbol{\lambda} \in (1-\epsilon)\Lambda^*$ for
some $\epsilon > 0$), the BPE max-weight payment routing (\Cref{def:maxweight})
stabilizes all virtual queues. Specifically:
\begin{equation}
\limsup_{T \to \infty} \frac{1}{T} \sum_{t=0}^{T-1} \sum_{k,\tasktype}
\E\bigl[Q(k, t, \tasktype)\bigr] \leq \frac{D}{\epsilon}
\end{equation}
where $D$ is a constant depending on maximum arrival and capacity rates.
\end{theorem}

\begin{proof}[Proof sketch]
Following Neely~\cite[Ch.~4]{neely2010stochastic}, we bound the Lyapunov drift.
From the queue dynamics~\eqref{eq:queue}:
\begin{align}
Q(k, t+1, \tasktype)^2 &\leq Q(k, t, \tasktype)^2
    + F(k, t, \tasktype)^2 + \Csmooth(k, t, \tasktype)^2 \notag \\
    &\quad + 2Q(k, t, \tasktype)\bigl[F(k, t, \tasktype) - \Csmooth(k, t, \tasktype)\bigr]
\end{align}

Summing over all $(k, \tasktype)$ and taking conditional expectations:
\begin{equation}
\Delta(t) \leq D + 2 \sum_{k, \tasktype} Q(k, t, \tasktype) \cdot
\E\bigl[F(k, t, \tasktype) - \Csmooth(k, t, \tasktype) \;\big|\; \mathbf{Q}(t)\bigr]
\end{equation}
where $D = \sum_{k,\tasktype} \bigl[F_{\max}^2 + C_{\max}^2\bigr]$.

The max-weight rule~\eqref{eq:maxweight} minimizes the right-hand side over all
feasible allocations. Since $\boldsymbol{\lambda} \in (1-\epsilon)\Lambda^*$,
there exists a stationary randomized policy with $\E[F(k,t,\tasktype)] \leq
\Csmooth(k,\tasktype) - \epsilon'$ for some $\epsilon' > 0$. The max-weight
policy achieves drift at least as negative, yielding:
\begin{equation}
\Delta(t) \leq D - \epsilon' \sum_{k,\tasktype} Q(k, t, \tasktype)
\end{equation}

By the Foster--Lyapunov criterion~\cite{meyn2009markov}, the queue process
is positive recurrent, and the time-averaged backlog bound follows from
telescoping the drift inequality.
\end{proof}

\subsection{Overflow Buffer Bound}

\begin{theorem}[Bounded Buffer]
\label{thm:buffer}
Under the same conditions as \Cref{thm:main}, the overflow buffer satisfies:
\begin{equation}
\limsup_{T \to \infty} \frac{1}{T} \sum_{t=0}^{T-1}
\E\bigl[B(\tasktype, t)\bigr] \leq \frac{D_B}{\epsilon}
\end{equation}
for a constant $D_B$ depending on the maximum arrival-capacity gap.
\end{theorem}

\begin{proof}[Proof sketch]
Augment the Lyapunov function with the buffer: $\tilde{L}(t) = L(t) + \sum_\tasktype B(\tasktype, t)^2$.
The buffer dynamics mirror a virtual queue with service rate equal to total
available capacity. The drift bound carries through identically.
\end{proof}

\begin{remark}
The no-drop constraint distinguishes BPE from standard backpressure. In data
networks, excess packets are dropped; here excess payments are buffered. The
buffer bound guarantees that this does not cause unbounded accumulation under
sufficient aggregate capacity.
\end{remark}
